\section{A terminating prescription for the witnesses}
\label{sec:prescription}

We now fix every choice in the preceding construction. The purpose of the prescription is to define one field, one group, and two finite sums without relying on an unspecified successful random choice.

\subsection{Finite search for the incidence data}

Order binary coefficient vectors lexicographically, with $0<1$; for a polynomial of fixed degree use the list of coefficients in increasing degree order. For each integer $h\ge2$, perform the following finite computations.
\begin{enumerate}[label=\textup{(\arabic*)}]
\item Put $q=2^h$, and let $\ell$ be the least prime congruent to $-1$ modulo $q^{100}$. Test successive integers in that class for primality. Linnik's theorem guarantees that this search terminates. Define $D=\ell+1$, $s=\lfloor\log_qD\rfloor$, $b=s+20$, and $n=10b$ as in Section~\ref{sec:hypergraph}.
\item Let $g$ be the lexicographically first monic irreducible polynomial over $\F_2$ of degree $h(\ell-1)$, and put
\begin{equation}
\label{eq:exact-field}
K=\F_2[Z]/(g(Z)).
\end{equation}
Order $K$ by the coefficient vectors of its unique representatives of degree less than $h(\ell-1)$, padded by zeros. Irreducible polynomials of this degree exist, and irreducibility is a finite test.
\item Take the subfield
\[
Q=\{x\in K:x^q=x\},
\]
with its induced order. It has $q$ elements since $h$ divides $h(\ell-1)$. Let $\zeta$ be the first element of $K$ of multiplicative order $\ell$. Such an element exists: Fermat's little theorem gives $\ell\mid 2^{\ell-1}-1$, so $\ell$ divides $|K^\times|=2^{h(\ell-1)}-1$, and $K^\times$ is cyclic.
\item Form $\mathcal A=Q^n$ in lexicographic order and the index simplex $S$ of Section~\ref{sec:pairing}. Order subsets of a finite ordered set by their indicator vectors, with $0<1$. In particular this orders affine lines as subsets of $\mathcal A$, families of lines as subsets of the ordered line set, and pairs $(O,\mathcal F)$ lexicographically. Test every pair with $O\subset S$ and $\mathcal F$ a family of full affine $Q$-lines for the following conditions:
\begin{enumerate}[label=\textup{(\alph*)}]
\item $2q^n/D^2\le |O|\le7q^n/D^2$;
\item every line of $\mathcal F$ is contained in $\mathcal A\setminus O$;
\item every point of $\mathcal A\setminus O$ lies on exactly $D$ lines of $\mathcal F$;
\item the incidence graph has no cycle of length less than $12$.
\end{enumerate}
\end{enumerate}
Start with $h=2$. For a fixed $h$, step~\textup{(4)} tests only finite sets and a finite graph. Stop at the first successful pair; if every pair fails, increase $h$ by one and repeat. This selects the least successful $h$ and then the first successful pair in its prescribed order. Theorem~\ref{thm:hypergraph} guarantees success for every sufficiently large $h$, uniformly in the enumeration of $Q$, so the outer search terminates. The search enumerates the required incidence data directly, not the auxiliary random choices used to prove their existence. The auxiliary extension field in that proof is only a model for a $Q$-vector space; it is not required to embed in $K$.

Fix these data and write $m=|O|$, $t=|V|$, where $V=\mathcal A\setminus O$. The indicator polynomials $f_o$ and vectors $x_v$ are the unique ones from Proposition~\ref{prop:pairing}; their coefficients may be computed by finite-field linear algebra. We use the induced orders on $O$, $V$, and $\mathcal F$.

\subsection{Local groups, the presentation, and the two sums}

At each point $v$, use the standard coordinate copy $K_v^0=(\F_\ell^2,+)$, and for each line $f$ use $L_f^0=(\F_\ell,+)$. Order $\F_\ell$ by the residues $0,\ldots,\ell-1$ and $\F_\ell^2$ lexicographically. Order its one-dimensional subspaces as subsets. Match them in order to the $D=\ell+1$ incident lines at $v$, and define $\iota_{vf}$ by sending $1\in L_f^0$ to the first nonzero vector of the assigned direction. These maps specify the group $N$ in Lemma~\ref{lem:embedding} without any further choices.

Choose $v_*$ to be the first point of $V$. The isomorphism $\phi_v:K_v\to K_{v_*}$ preserves the standard coordinates transported from the embedded copies. Order linear forms on $\F_\ell^2$ by their coefficient pairs. Assign the points of $V$, in order, to the first $t$ pairs
\[
(s,\alpha),\qquad 1\le s\le m,\quad \alpha\in(\F_\ell^2)^*,
\]
in lexicographic order. Proposition~\ref{prop:pairing} gives $t\le m\ell^2$, so there are enough pairs. This fixes every slot and every character in the construction of Theorem~\ref{thm:criterion}.

For the finite presentation of $N$, use the local zero-sum relators of length at most three, as justified in Section~\ref{sec:criterion}. Section~\ref{sec:packing} then defines $H=N\times\langle z\rangle$, the conjugation extension $L$, the finite group $F=(\mathbb Z/\ell)^m\rtimes S_m$, and finally $G=L\times F$; Section~\ref{sec:finite-presentation} gives the resulting finite presentation. The regular-permutation action used to prove that $H$ embeds in $L$ is an embedding certificate; a choice of such permutations is not part of the definition of the group.

With this group fixed, the formulas in Section~\ref{sec:packing} give $e,P,U,T,C,J_0$, the two square matrices $a_{\rm mat},b_{\rm mat}$, and the matrix units $p_{ij}$. The two required finite expressions are
\begin{equation}
\label{eq:exact-outputs}
\begin{split}
a_{\mathrm{out}}&=1-\sum_{i=1}^m p_{ii}
 +\sum_{i,j=1}^m(a_{\rm mat})_{ij}p_{ij},\\
b_{\mathrm{out}}&=1-\sum_{i=1}^m p_{ii}
 +\sum_{i,j=1}^m(b_{\rm mat})_{ij}p_{ij}.
\end{split}
\end{equation}
Coefficients from $K[L]$ and matrix units from $K[F]$ are viewed in the two factors of $K[G]$. The displayed expressions, together with the preceding formulas, involve only finitely many group words and field coefficients. They define finite sums whether or not distinct displayed words represent the same group element; no word-equality oracle is used to normalize them.

\begin{proof}[Proof of Theorem~\ref{thm:main}]
The finite search terminates by Theorem~\ref{thm:hypergraph}. Its output satisfies Proposition~\ref{prop:pairing}, and hence every hypothesis of Theorem~\ref{thm:criterion}. That theorem proves the two product assertions for \eqref{eq:exact-outputs}, supplies a finite presentation of this group $G$, and retains an embedded element of order $\ell$. In particular the nonvanishing of the reverse-product defect follows by specializing the right-hand side of Equation~\eqref{eq:reverse-sandwich}, which lies in $\Mat_t(K[H])$, to a nonzero matrix over $K$; it does not depend on an unproved inequality between group words. The prime $\ell$ is odd, and the field in \eqref{eq:exact-field} is finite of characteristic two.
\end{proof}

\subsection{Scope of the construction}

The construction uses torsion explicitly. The embedded subgroups of orders $\ell$, $\ell^2$, and $\ell^3$ have odd order, so every denominator in their averages is invertible in $K$. In the finite auxiliary group we invert only $|(\mathbb Z/\ell)^m|$, never $|S_m|$. Characteristic two makes the regular degree $\ell+1$ vanish as a scalar, allowing the simultaneous line and total pairing identities of Proposition~\ref{prop:pairing}. No assertion about group algebras of torsion-free groups is needed.

For comparison, if $G_0$ is finite, then $K_0[G_0]$ is finite dimensional over every field $K_0$, even when the characteristic divides $|G_0|$. Left multiplication is a faithful representation of this algebra. In finite dimension, $L_aL_b=I$ implies $L_bL_a=I$, and hence $ab=1$ implies $ba=1$. Thus the group in Theorem~\ref{thm:main} is necessarily infinite.

Finally, for any group and any two finite group-algebra sums, the union of their supports generates a finitely generated subgroup. Its group algebra embeds in the original one, so both product assertions already hold or fail in that subgroup algebra. Finite generation is therefore automatic for a direct-finiteness counterexample after passing to the support subgroup. The stronger finite-presentation conclusion for the particular group $G$ here comes from the presentation argument in Section~\ref{sec:finite-presentation}; it does not follow merely from finite generation of the support subgroup.
